\section{Descent with compact isometric monodromy}
\label{V:isometric-descent}

We first apply the invariant results when the residual monodromy lies
in a compact torus. The geometric structure is supplied by the
finite-cover theorem of \emph{Anticanonical nonvanishing from smooth semipositivity} \cite[Theorem~2.1]{CompanionH}, based
on Demailly--Peternell--Schneider and Campana--Demailly--Peternell
\cite{DPS1996,CDP2015}. Its nonvanishing theorem is not used in the
argument below.

Let $X$ be smooth connected projective with smoothly semipositive
$-K_X$. The cited structure theorem gives a connected finite \'etale
cover $\nu:S\to X$ and an equivariant product description of its
universal cover
\[
 \widetilde X=\C^a\times F\times Z,
 \qquad S=(\C^a\times F\times Z)/\Lambda.
\]
Here $\Lambda$ acts by translations on $\C^a$, trivially on the
compact simply connected Ricci-flat projective factor $F$, and through
a compact connected real torus on $Z$. The variety $Z$ is smooth
projective, its anticanonical line is smoothly semipositive, and
$H^i(Z,\OO_Z)=0$ for every $i>0$. The compact action has algebraic
torus closure $G\subset\operatorname{Aut}(Z)$, with its natural
anticanonical linearization. There are invariant canonical frames
$\eta_0,\eta_F$ on $\C^a,F$. Any factor may be a point.

Theorem~\ref{V:invariant-main} gives a nonzero
$s\in H^0(Z,K_Z^{-m})^G$ for some $m>0$. If $Z$ is a point, take
$m=1$ and $s=1$. Equally, the controlled-field conversion and its
index application can be used here: $Z$ has irregularity zero and
actual index one, and its natural compact action meets precisely those
hypotheses.

Under the determinant identity for the product tangent bundle, the
section
\[
 \operatorname{pr}_0^*(\eta_0^{-m})\otimes
 \operatorname{pr}_F^*(\eta_F^{-m})\otimes
 \operatorname{pr}_Z^*s
\]
belongs to $H^0(\widetilde X,K_{\widetilde X}^{-m})$. It is nonzero,
since the first two factors never vanish and $s$ is nonzero on a
nonempty open set of $Z$. Translations preserve $\eta_0$, the action
on $F$ preserves $\eta_F$, and the residual action on $Z$ fixes $s$.
Thus this product is $\Lambda$-invariant and descends to a nonzero
holomorphic section of $K_S^{-m}$. It is algebraic by projectivity.
The finite \'etale norm \cite[Lemma~2.2]{CompanionH}
then gives a nonzero section of $K_X^{-m\deg\nu}$.

The product is used on the universal cover; the finite cover $S$ need
not split as a product. The next argument starts from a locally constant
rationally connected fibration instead. Its algebraic monodromy can
have a nontrivial unipotent part, which requires an additional
fixed-vector argument.
